\documentclass[11pt]{article}
\usepackage{amsmath}
\usepackage[hidelinks]{hyperref}
\begin{document}
\title{Return-to-contact triggers and an ACK-retained old duty}
\author{}
\date{26 November 2026}
\maketitle

\paragraph{Audit and result.}
A queue still containing a report can represent two distinct
operational states: base already applied it, but the return
acknowledgment was lost; or the earlier receive reached a
terminally unserved application state and must be recovered.
In the modeled mission those states can have the same queue
size, ACK history, map, and offered link rates. A timely, scoped
terminal receipt lets the mission avoid one irrevocable staging
commitment in the served case and still guarantee the old report,
new inspection, and egress in the unserved case. The strongest
aggregate-selected menu feasible in both states always stages.
The exact direct decision rule with the same typed receipt ties
the shared checker. The gain is conditional and pointwise against
the specified aggregate movement interfaces, not against all
implementations of ACHORD or an adapted scheduler.

\paragraph{Instance and source cut.}
At tick 0 the original packet reaches the base endpoint in both
states. The report owner either pays the obligation at tick 1
or closes that attempt without payment. No attempt remains in
flight on that application path. No acknowledgment clears the sender
queue before the first choice, so packet delivery and outstanding
service are not interchangeable. The base handler has idempotent
completion for the same report identity: an unnecessary transport
copy costs bytes but does not pay the same obligation twice.
Application lifetime, idempotence and ownership are implemented in
the test harness around native packet callbacks; they are not
properties inferred from the SubT broker.

The current lead's cut is tick 3. A one-byte message from the
licensed terminal source at tick 2 may report the current result, be
erased, or refer to older authority. A delayed message at tick 4
cannot certify the earlier cut. Receiver-visible channel output,
not the raw source branch, selects the plan at the cut. Both
compared mechanisms see the same signal and pay the same baseline
read attempt; the literal movement trigger does not use terminal
semantics. An old report arriving after the modeled closed
attempt, different report IDs, misleading valid payloads, and
source authority changes require a distinct experiment.

\paragraph{Physical and radio envelope.}
A declared known directed civil inspection graph starts at a
junction J. No primary link at J can pay the old application
duty. Direct continuation goes J--I--E for inspection and egress
at ticks 4 and 7, giving an old retry at 7. Return via the strong
comms checkpoint C gives an old retry at 5, but inspection at
8 and egress at 11. The hard deadlines are report 6, inspection
6, and egress 10. A third route J--M--I--E engages one
non-reusable physical staging slot, then makes inspection at 6
and egress at 9. There is no path that can start this commitment
one tick later and still meet both deadlines on the unserved
branch. The auxiliary relay's modeled egress is at tick 8;
an independent actual callback at 8 is also required and survives
all the tested plans. This pilot does not infer field egress
paths or learn an unknown map.

The staged mule uses one route, not two simultaneous reserves.
A local read at the mule at tick 4 selects contact A or B, can
fail to arrive, and costs one transmitted byte when staging
was engaged. Even with that read, both contacts can be lost
outside the positive envelope. If A becomes unavailable before
settlement, only a received local switch report enables use of
B. Two-hop report transport at 5 and 6, the independently routed
work, the terminal receipt, and the read all exercise the pinned
SubT broker/RF code. The graph, fault law of the link, and
clock/space association are supplied to the harness, not found
by the broker. A staged option can suppress duplicate transport
using the terminal state available at its receiving handler;
this same generous late suppression is granted to the comparator.
It does not refund the slot or the local read.

\paragraph{Compared mechanisms.}
In \href{https://lillyclark.github.io/files/ACHORD_RA_L_IROS_.pdf}{Saboia et al.,
ACHORD}, sections III-B and IV, topic buffers are ACK-retained,
with reliability and traffic classes and visible buffer/rate/estimated
drain metrics. The published return-to-comms behavior enters above
300 KB by default, heads to a strong nearby checkpoint, resumes
below 200 KB, and can move toward the base after sustained lack
of drain. The paper also uses deployment coordination and discusses
a successful robot relay. The executable comparison here is the
specified motion trigger on an authored graph with one reachable
strong return checkpoint; it does not reimplement the entire
network or exploration pipeline, and no timeout occurs before the
irrevocable choice. The actual paired source-stage results should
not be attributed to a tunable high-water mark. Threshold and
hysteresis variations cannot distinguish inputs with identical
backlog and prior mode. Returning or continuing fails a hard
deadline on the unserved branch.

To give the aggregate architecture a stronger return, let it
choose any route, stage or not, and suppress redundant later sends.
Every robust first choice in that menu still stages, because
otherwise the unserved case cannot meet both deadlines. A
new controller evaluates whole history records and chooses the
least costly admitted option using the current scoped receipt.
This preserves the same mission safety and all prescribed service.
It uses zero instead of one staging commitment on the served
branch, no additional payload traffic, and reaches inspection
and exit two ticks earlier; the unserved branch is a tie.
Worst-case required stages remains one. No expected gain is
claimed without a distribution. If the scoped receipt is absent,
old or late, the unchanged checker also must stage to serve both;
loss of the switch read or both backup links can instead make
the unserved class unachievable in the catalogue. A direct
factored policy given the same roles, signal and choices gives
the exact same result.

\paragraph{Checks and reserve.}
The development replay checks 160 native cases, 320 full
checker/direct comparisons, 39,852 valid threshold settings
with two prior modes over six illustrative backlogs, and four
actual inline native selections after observing the earlier
signal. The finite threshold sweep is a test; the impossibility
for a fixed aggregate trigger follows separately because it
must make the same choice on the paired inputs. Lean checks the
minimum-stage argument in the declared positive menu. No
evidentiary claim is made about a full ACHORD or HeRo battery,
an independent mission, or the architecture acceptance
disposition. A stronger study would include learned reachability
and an independently frozen comparator capable of joint return
and rendezvous as in
\href{https://ntrs.nasa.gov/citations/20220001447}{Klaesson et al.};
here the joint timed source and receiver constraints are a
design-case illustration rather than a substitute for it.
\end{document}
